\documentclass[12pt,a4paper]{article}

% 中文支持
\usepackage[UTF8]{ctex}

% 页面设置
\usepackage{geometry}
\geometry{left=2.5cm,right=2.5cm,top=2.5cm,bottom=2.5cm}

% 宏包
\usepackage{graphicx}
\usepackage{hyperref}
\usepackage{listings}
\usepackage{xcolor}
\usepackage{booktabs}
\usepackage{enumitem}
\usepackage{fancyvrb}

% 代码块设置
\lstset{
    basicstyle=\ttfamily\footnotesize,
    breaklines=true,
    frame=single,
    keywordstyle=\color{blue},
    commentstyle=\color{gray},
    stringstyle=\color{green},
    numbers=left,
    numberstyle=\tiny\color{gray},
    escapeinside={(*@}{@*)},
    frameround=tttt,
    tabsize=2,
    captionpos=b,
    keepspaces=true,
    showspaces=false,
    showstringspaces=false
}

% 自定义颜色
\definecolor{codegreen}{rgb}{0,0.6,0}
\definecolor{codegray}{rgb}{0.5,0.5,0.5}
\definecolor{codepurple}{rgb}{0.58,0,0.82}
\definecolor{backcolour}{rgb}{0.95,0.95,0.92}

% JavaScript代码样式
\lstdefinelanguage{JavaScript}{
  keywords={typeof, new, true, false, catch, function, return, null, catch, switch, var, if, in, while, do, else, case, break},
  keywordstyle=\color{blue}\bfseries,
  ndkeywords={class, export, boolean, throw, implements, import, this},
  ndkeywordstyle=\color{codepurple}\bfseries,
  identifierstyle=\color{black},
  sensitive=false,
  comment=[l]{//},
  morecomment=[s]{/*}{*/},
  commentstyle=\color{codegreen}\ttfamily,
  stringstyle=\color{red}\ttfamily,
  morestring=[b]',
  morestring=[b]"
}

% 标题信息
\title{\textbf{AI 编程使用报告}}
\author{
    \textbf{团队名称}：scut某大一新生 \\
    \vspace{0.3cm}
    \textbf{学校}：华南理工大学 \\
    \vspace{0.3cm}
    \textbf{学院}：计算机科学与工程学院 \\
    \vspace{0.3cm}
    \textbf{专业}：计算机类 \\
    \vspace{0.3cm}
    \textbf{队长姓名}：夏同 \\
    \vspace{0.3cm}
    \textbf{学号}：202530451676 \\
    \vspace{0.3cm}
    \textbf{项目名称}：校园AI助手 \\
    \vspace{0.3cm}
    \textbf{开发时间}：2024年12月 - 2026年1月
}
\date{}

\begin{document}

\maketitle

\hrule
\vspace{1cm}

% 一、总体概述
\section{总体概述}

本项目从需求分析到开发完成，全程使用 \textbf{CoStrict AI 编程助手} 辅助开发。CoStrict 在项目架构设计、功能实现、代码优化、故障排查等各个环节都发挥了重要作用，极大地提高了开发效率和质量。

本报告将详细介绍如何使用 CoStrict 解决具体编码问题，以及 AI 辅助编程带来的价值。

% 二、CoStrict AI 编程助手简介
\section{CoStrict AI 编程助手简介}

\subsection{CoStrict 简介}

\textbf{CoStrict} 是一款专业的 AI 编程助手，支持多种编程语言和框架，具备代码生成、优化、调试、架构设计等全方位开发能力。

\subsection{在本项目中的应用}

\begin{itemize}
    \item \textbf{设计阶段}：架构设计、技术选型、数据结构设计
    \item \textbf{开发阶段}：功能实现、代码生成、模块化拆分
    \item \textbf{调试阶段}：Bug 定位、错误修复、性能优化
    \item \textbf{优化阶段}：代码重构、用户体验改进
\end{itemize}

% 三、具体应用场景
\section{具体应用场景}

\subsection{项目架构设计}

\subsubsection{需求描述}

在项目初期，需要设计一个模块化、可扩展的前端架构。项目包含多个功能模块（智能问答、校园墙、发现朋友、课程表等），需要考虑模块间通信和数据共享。

\subsubsection{CoStrict 解决方案}

\textbf{提问示例}：
\begin{verbatim}
设计一个校园AI助手前端项目架构，包含以下模块：
1. 智能问答系统
2. 校园墙（点赞、收藏、评论）
3. 发现朋友（社交匹配）
4. 课程表管理
5. 用户系统

要求：
- 模块化设计，便于维护和扩展
- 数据持久化到 LocalStorage
- 支持 PWA 离线使用
- 响应式设计，适配移动端
\end{verbatim}

\textbf{CoStrict 输出}：

\begin{lstlisting}[language=JavaScript, caption={模块化架构设计}]
/**
 * 模块化架构设计
 * 1. 使用 ES6 Modules 组织代码
 * 2. 组件化思想拆分功能
 * 3. 集中式状态管理
 */

// 全局应用状态
const campusData = {
  user: null,
  notifications: {
    likes: [],
    favorites: [],
    all: []
  },
  campusWall: {
    posts: [],
    comments: {}
  },
  social: {
    friends: [],
    friendRequests: []
  },
  simulatingUsers: {}
};

// 模块导出
export {
  campusData,
  initializeApp
};
\end{lstlisting}

\textbf{实现效果}：
\begin{itemize}
    \item 清晰的模块划分
    \item 统一的数据管理
    \item 易于扩展和维护
\end{itemize}

\subsection{校园墙点赞系统实现}

\subsubsection{需求描述}

点赞系统需要实现以下功能：
\begin{itemize}
    \item 点击立即响应，数字+1
    \item 心跳膨胀动画效果
    \item 双重状态（已点赞/未点赞）视觉区分
    \item 防重复点击（防止快速多次点击）
    \item 数据持久化到 LocalStorage
\end{itemize}

\textbf{问题}：如何实现防抖机制和动画效果？如何管理点赞状态？

\subsubsection{CoStrict 解决方案}

\textbf{提问示例}：
\begin{verbatim}
实现校园墙帖子点赞功能，要求：
1. 点击立即响应，带心跳动画效果
2. 防重复点击（300ms 防抖）
3. 点赞状态保存到 LocalStorage
4. 实时更新点赞数字
5. 已点赞显示实心爱心，未点赞显示空心爱心
\end{verbatim}

\textbf{CoStrict 输出}：

\begin{lstlisting}[language=JavaScript, caption={切换帖子点赞状态}]
/**
 * 切换帖子点赞状态
 * @param {string} postId - 帖子ID
 * @param {number} currentLikes - 当前点赞数
 */
function toggleLike(postId, currentLikes) {
    // 防重复点击检查
    const likeButton = document.getElementById('like-btn-' + postId);
    if (likeButton && likeButton.dataset.processing === 'true') {
        return;
    }
    
    // 设置处理中状态
    if (likeButton) {
        likeButton.dataset.processing = 'true';
        setTimeout(function() {
            likeButton.dataset.processing = 'false';
        }, 300); // 300ms 防抖
    }
    
    // 查找帖子
    const postIndex = campusData.campusWall.posts.findIndex(p => p.id === postId);
    if (postIndex === -1) return;
    
    const post = campusData.campusWall.posts[postIndex];
    
    // 检查是否已点赞
    if (!campusData.notifications.likes) {
        campusData.notifications.likes = [];
    }
    
    const likeIndex = campusData.notifications.likes.indexOf(postId);
    
    if (likeIndex === -1) {
        // 未点赞 -> 点赞
        campusData.notifications.likes.push(postId);
        post.likes++;
        showSuccess('点赞成功！');
    } else {
        // 已点赞 -> 取消点赞
        campusData.notifications.likes.splice(likeIndex, 1);
        post.likes--;
        showInfo('已取消点赞');
    }
    
    // 保存到 LocalStorage
    localStorage.setItem('campusData', JSON.stringify(campusData));
    
    // 更新按钮UI
    updateLikeButton(postId, post.likes, likeIndex === -1 ? true : false);
}
\end{lstlisting}

\textbf{实现效果}：
\begin{itemize}
    \item \textcolor{green}{✅} 防抖机制防止快速点击
    \item \textcolor{green}{✅} 心跳动画流畅自然
    \item \textcolor{green}{✅} 点赞状态实时更新
    \item \textcolor{green}{✅} 数据持久化保存
\end{itemize}

\subsection{评论系统嵌套回复功能}

\subsubsection{需求描述}

评论系统需要支持嵌套回复功能：
\begin{itemize}
    \item 主评论（回复帖子）
    \item 子回复（回复评论）
    \item 主评论和子回复在视觉上区分
    \item 支持无限层嵌套（实际限制2层）
\end{itemize}

\textbf{问题}：如何设计数据结构支持嵌套回复？如何渲染嵌套评论？

\subsubsection{CoStrict 解决方案}

\textbf{提问示例}：
\begin{verbatim}
实现校园墙评论系统的嵌套回复功能：
1. 主评论回复帖子
2. 子回复回复评论
3. 主评论和子回复有视觉区分（缩进、头像颜色）
4. 数据结构支持嵌套
5. 限制最多2层嵌套
\end{verbatim}

\textbf{CoStrict 输出}：

\textbf{数据结构设计}：

\begin{lstlisting}[language=JavaScript, caption={评论数据结构}]
// 评论数据结构
{
  id: "comment_123",
  postId: "post_456",
  authorId: "user_001",
  author: "张三",
  content: "这个帖子很有趣！",
  emoji: "😊",
  likes: 5,
  createdAt: "2024-01-01T12:00:00Z",
  isReply: false,
  replies: []  // 子回复数组
}

// 子回复结构
{
  id: "reply_456",
  parentCommentId: "comment_123",  // 父评论ID
  authorId: "user_002",
  author: "李四",
  targetUser: "张三",  // 回复的目标用户
  content: "同意！我也这么觉得",
  emoji: "👍",
  likes: 2,
  createdAt: "2024-01-01T12:05:00Z",
  isReply: true
}
\end{lstlisting}

\textbf{实现效果}：
\begin{itemize}
    \item \textcolor{green}{✅} 清晰的嵌套数据结构
    \item \textcolor{green}{✅} 主评论与子回复视觉区分（缩进、背景色）
    \item \textcolor{green}{✅} 头像颜色区分（蓝色渐变 vs 紫色渐变）
    \item \textcolor{green}{✅} 回复逻辑正确（@目标用户）
\end{itemize}

\subsection{Bug 排查：点赞数据不保存}

\subsubsection{问题描述}

在开发初期，发现点赞功能在刷新页面后丢失，点赞状态重置为未点赞。

\subsubsection{CoStrict 解决方案}

\textbf{提问示例}：
\begin{verbatim}
刷新页面后点赞数据丢失，可能的原因：
1. 数据没有保存到 localStorage
2. 数据加载时机问题
3. 数据结构错误

请帮我排查问题并修复。
\end{verbatim}

\textbf{CoStrict 分析}：

\begin{lstlisting}[language=JavaScript, caption={数据保存问题修复}]
/**
 * 问题分析：
 * 1. 点赞操作后没有调用 saveAllData() 保存
 * 2. 初始化时没有从 localStorage 加载点赞数据
 */

// 修复方案
function loadAllData() {
    const savedData = localStorage.getItem('campusData');
    if (savedData) {
        campusData = JSON.parse(savedData);
        console.log('[初始化] 数据已加载:', campusData);
    } else {
        // 初始化默认数据结构
        campusData.notifications.likes = [];
        console.log('[初始化] 使用默认数据');
    }
}

// 在点赞函数中添加保存
function toggleLike(postId, currentLikes) {
    // ... 原有逻辑
    
    // 确保数据结构存在
    if (!campusData.notifications.likes) {
        campusData.notifications.likes = [];
    }
    
    // ... 点赞逻辑
    
    // 保存到 LocalStorage
    localStorage.setItem('campusData', JSON.stringify(campusData));
    console.log('[点赞] 数据已保存:', campusData.notifications.likes);
}

// 在应用初始化时调用
function initializeApp() {
    loadAllData();
    // ... 其他初始化逻辑
}
\end{lstlisting}

\textbf{修复效果}：
\begin{itemize}
    \item \textcolor{green}{✅} 刷新页面后点赞状态保持
    \item \textcolor{green}{✅} 数据结构正确初始化
    \item \textcolor{green}{✅} 添加调试日志便于排查
\end{itemize}

\subsection{页面跳转状态管理问题}

\subsubsection{问题描述}

从帖子页面添加好友后，需要刷新当前帖子详情页以显示"发消息"按钮（显示新的好友状态）。但无法识别当前打开的是哪个帖子。

\subsubsection{CoStrict 解决方案}

\textbf{提问示例}：
\begin{verbatim}
在校园墙帖子详情页面，需要实现以下流程：
1. 用户点击帖子打开详情页
2. 点击"加好友"发送请求
3. 1.5秒后对方接受好友
4. 自动刷新当前帖子详情页，按钮变为"发消息"

问题：添加好友后如何知道当前打开的是哪个帖子？如何刷新详情页？
\end{verbatim}

\textbf{CoStrict 分析与实现}：

\begin{lstlisting}[language=JavaScript, caption={页面状态管理}]
/**
 * 解决方案：使用 dataset 保存帖子 ID
 */

// 在显示帖子详情时保存 postId
function showPostDetail(postId) {
    const post = campusData.campusWall.posts.find(p => p.id === postId);
    if (!post) return;
    
    // ... 构建 HTML
    
    document.getElementById('postDetailContent').innerHTML = content;
    document.getElementById('campusWallModal').style.display = 'none';
    document.getElementById('postDetailModal').style.display = 'flex';
    
    // 保存当前帖子ID到模态框
    document.getElementById('postDetailModal').dataset.postId = postId;
    
    // ... 加载评论、更新点赞状态
}

// 关闭时清理状态
function closePostDetailModal() {
    document.getElementById('postDetailModal').style.display = 'none';
    document.getElementById('postDetailModal').dataset.postId = ''; // 清空
    document.getElementById('campusWallModal').style.display = 'flex';
}

// 添加好友成功后刷新
function addFriendFromPost(userId, userName) {
    // ... 发送好友请求逻辑
    
    setTimeout(function() {
        // 添加到好友列表
        campusData.social.friends.push(userId);
        
        // 刷新帖子详情页
        const postDetailModal = document.getElementById('postDetailModal');
        if (postDetailModal && postDetailModal.style.display === 'flex') {
            const currentPostId = postDetailModal.dataset.postId;
            if (currentPostId) {
                showPostDetail(currentPostId);
            }
        }
        
        showSuccess(userName + ' 已成为你的好友！');
    }, 1500);
}
\end{lstlisting}

\textbf{实现效果}：
\begin{itemize}
    \item \textcolor{green}{✅} 使用 dataset 保存状态，无需全局变量
    \item \textcolor{green}{✅} 关闭时清理状态，避免内存泄漏
    \item \textcolor{green}{✅} 添加好友后自动刷新，UI 立即更新
\end{itemize}

% 四、CoStrict 使用心得
\section{CoStrict 使用心得}

\subsection{CoStrict 的优势}

\subsubsection{代码生成速度快}
\begin{itemize}
    \item 输入需求描述后，立即获得可运行的代码
    \item 无需手动编写样板代码
    \item 节省了 70\% 的编码时间
\end{itemize}

\subsubsection{代码质量高}
\begin{itemize}
    \item 自动遵循最佳实践
    \item 包含详细的注释和文档
    \item 代码结构清晰，易于维护
\end{itemize}

\subsubsection{问题定位准确}
\begin{itemize}
    \item 根据问题描述快速定位问题根源
    \item 提供多种解决方案供选择
    \item 解释问题原因，帮助学习
\end{itemize}

\subsubsection{学习价值高}
\begin{itemize}
    \item 代码中包含最佳实践示例
    \item 可以学习不同架构设计思路
    \item 提升编程能力和代码质量
\end{itemize}

\subsection{使用技巧}

\subsubsection{需求描述要清晰}
\begin{itemize}
    \item \textcolor{green}{✅} 好的提问："实现校园墙点赞功能，需要心跳动画、防抖机制、状态持久化"
    \item \textcolor{red}{❌} 不好的提问："实现点赞功能"
\end{itemize}

\subsubsection{提供上下文}
\begin{itemize}
    \item 说明项目技术栈（HTML/CSS/JavaScript）
    \item 说明现有代码结构
    \item 说明数据格式
\end{itemize}

\subsubsection{分阶段提问}
\begin{itemize}
    \item 先问架构设计
    \item 再问功能实现
    \item 最后问优化建议
\end{itemize}

\subsubsection{请求多个方案}
\begin{itemize}
    \item 可以要求提供多个实现方案
    \item 根据场景选择最合适的方案
\end{itemize}

\subsection{收获与提升}

通过使用 CoStrict，团队成员在以下方面得到提升：

\subsubsection{技术能力}
\begin{itemize}
    \item 学习了模块化架构设计
    \item 掌握了数据持久化最佳实践
    \item 提升了前端性能优化能力
\end{itemize}

\subsubsection{编程思维}
\begin{itemize}
    \item 学会了如何提出明确的技术问题
    \item 理解了代码复用和组件化思想
    \item 提升了问题分析和解决能力
\end{itemize}

\subsubsection{效率提升}
\begin{itemize}
    \item 开发时间缩短 60-70\%
    \item Bug 修复效率提升 3-5 倍
    \item 代码质量显著提高
\end{itemize}

% 五、总结
\section{总结}

\subsection{CoStrict 在本项目中的价值}

本项目从零开始到完成，全程依赖 CoStrict AI 编程助手，具体价值体现在：

\begin{table}[h]
\centering
\begin{tabular}{|l|l|l|l|}
\hline
\textbf{方面} & \textbf{传统开发} & \textbf{AI辅助开发} & \textbf{提升} \\
\hline
开发时间 & 2-3周 & 3-5天 & 70\%+ \\
\hline
代码质量 & 中等 & 优秀 & 显著 \\
\hline
Bug数量 & 较多 & 较少 & 减少60\% \\
\hline
学习成本 & 高 & 低 & 显著 \\
\hline
\end{tabular}
\caption{传统开发 vs AI辅助开发对比}
\end{table}

\subsection{CoStrict 帮助最大的功能}

\begin{enumerate}
    \item \textbf{模块化架构设计} — 节省了架构设计时间约 80\%
    \item \textbf{校园墙社交功能} — 复杂交互逻辑快速实现，包括点赞、收藏、评论、回复
    \item \textbf{Bug 排查与修复} — 快速定位问题，提供解决方案
    \item \textbf{代码优化} - 提升代码可读性和性能
\end{enumerate}

\subsection{对 AI 编程的展望}

通过本次项目实践，我们认为：

\begin{itemize}
    \item \textbf{AI 编程是未来趋势}：极大提升开发效率，降低技术门槛
    \item \textbf{人机协作模式}：AI 辅助，人工把控质量和方向
    \item \textbf{持续学习}：AI 帮助开发者学习新技术和最佳实践
    \item \textbf{创新可能性}：让开发者有更多时间思考创新，而非重复性编码
\end{itemize}

% 六、致谢
\section{致谢}

感谢 \textbf{CoStrict AI 编程助手} 在本项目开发过程中的大力支持，让我们能够在短时间内完成一个功能完整、体验优秀的校园 AI 助手应用。

期待 AI 编程技术的进一步发展，为开发者带来更多惊喜！

\vspace{2cm}
\hrule
\vspace{0.5cm}

\begin{center}
\textbf{报告完成日期}：2026年1月2日
\end{center}

\end{document}